% from ega4-2-fr.tex
\oldpage[IV-2]{75}

Si $\mathscr F$ est réduit en $x$, il y a un voisinage ouvert $U$ de $x$ tel
que $\mathscr F|U$ soit réduit (3.2.2). En raisonnant comme dans (4.6.16), on
voit que dire que $\mathscr F$ est géométriquement réduit au point $x$
équivaut à dire que $\mathscr F$ est réduit au point $x$ et que le
sous-préschéma fermé $Y$ défini par l'annulateur $\mathscr J$ de $\mathscr F$
est géométriquement réduit au point $x$; il en résulte qu'il y a un voisinage
ouvert $U$ de $x$ tel que $\mathscr F|U$ soit géométriquement réduit (4.6.11).
Les raisonnements de (4.6.16) et (4.6.18) montrent aussi que si $X$ est
localement de type fini sur $k$, dire que $\mathscr F$ est géométriquement
ponctuellement intègre au point $x$ équivaut à dire que $\mathscr F$ est
réduit au point $x$ et que le sous-préschéma $Y$ est géométriquement
ponctuellement intègre au point $x$.

\subsection{Multiplicités dans la décomposition primaire sur un préschéma
algébrique}
\label{subsection:IV.4.7-fr}

\begin{lemma}[4.7.1]
\label{IV.2.4.7.1-fr}
Soient $A$, $B$ deux anneaux locaux, $\mathfrak m$ l'idéal maximal de $A$,
$\rho:A\to B$ un homomorphisme local. On suppose que $B$ soit un $A$-module
plat, et $B/\mathfrak mB$ un $B$-module de longueur finie. Alors, pour qu'un
$A$-module $M$ soit de longueur finie, il faut et il suffit que $M_{(B)}$ soit
un $B$-module de longueur finie, et on a
\begin{equation}
\operatorname{long}_B(M_{(B)})
=\operatorname{long}_A(M)\cdot\operatorname{long}_B(B/\mathfrak mB).
\tag{4.7.1.1}
\label{IV.2.4.7.1.1-fr}
\end{equation}
\end{lemma}

\begin{proof}
C'est un cas particulier de (2.5.5.2).
\end{proof}


\begin{corollary}[4.7.2]
\label{IV.2.4.7.2-fr}
Soient $A$, $B$ deux anneaux noethériens, $\rho:A\to B$ un homomorphisme
d'anneaux tel que $B$ soit un $A$-module plat. Soient $\mathfrak q$ un idéal
premier minimal de $B$, $\mathfrak p=\rho^{-1}(\mathfrak q)$, $M$ un
$A$-module de type fini; alors $\mathfrak p$ est un idéal premier minimal de
$A$, $M_{\mathfrak p}$ est un $A_{\mathfrak p}$-module de longueur finie et
l'on a
\begin{equation}
\operatorname{long}_{B_{\mathfrak q}}((M_{(B)})_{\mathfrak q})
=\operatorname{long}_{A_{\mathfrak p}}(M_{\mathfrak p})
 \operatorname{long}_{B_{\mathfrak q}}
 (B_{\mathfrak q}/\mathfrak pB_{\mathfrak q}).
\tag{4.7.2.1}
\label{IV.2.4.7.2.1-fr}
\end{equation}
\end{corollary}

\begin{proof}
En effet, $B_{\mathfrak q}$ est un $A_{\mathfrak p}$-module plat
$(0_{\mathrm I},6.3.2)$ et l'homomorphisme
$A_{\mathfrak p}\to B_{\mathfrak q}$ est local, donc $B_{\mathfrak q}$ est
un $A_{\mathfrak p}$-module fidèlement plat $(0_{\mathrm I},6.6.2)$; le fait
que $\mathfrak p$ est minimal dans $A$ résulte de (2.3.5). En outre,
$A_{\mathfrak p}$ est alors un anneau artinien (Bourbaki, \emph{Alg. comm.},
chap.~IV, \S~2, n\textsuperscript{o}~5, prop.~9) et $M_{\mathfrak p}$ est donc
un $A_{\mathfrak p}$-module de longueur finie; la formule (4.7.2.1) est donc
un cas particulier de (4.7.1.1).
\end{proof}

\begin{proposition}[4.7.3]
\EGAReaderTarget{ega4.c04.prop.4-7-3}{IV.2.4.7.3-fr}
Soient $k$ un corps d'exposant caractéristique $p$, $X$ un préschéma intègre
localement noethérien, $K=\mathrm R(X)$ le corps des fonctions rationnelles
sur $X$. Soient $k'$ une extension de $k$, $(X'_\alpha)$ la famille des
composantes irréductibles de $X'=X_{(k')}$, $x'_\alpha$ le point générique de
$X'_\alpha$.
\begin{enumerate}
\item[(i)] On suppose que $X'$ soit localement noethérien (ce qui a lieu si
$X$ est localement de type fini sur $k$, ou si $k'$ est une extension de type
fini de $k$). Soient alors $k''$ une extension quelconque de $k'$,
$X''=X\otimes_k k''$, $(X''_\beta)$ la famille des composantes irréductibles
de $X''$, $x''_\beta$ le point générique de $X''_\beta$. Si $x''_\beta$ est
au-dessus de $x'_\alpha$, on a
\begin{equation}
\operatorname{long}(\mathcal O_{x''_\beta})
=\operatorname{long}(\mathcal O_{x'_\alpha})
 \operatorname{long}(\mathcal O_{Z'',x''_\beta})
\tag{4.7.3.1}
\label{IV.2.4.7.3.1-fr}
\end{equation}
où l'on a posé $Z''=(X'_{\mathrm{red}})\otimes_{k'}k''$. En particulier, si
$k''$ est une extension séparable de $k'$ telle que $Z''$ soit localement
noethérien, on a
\[
\operatorname{long}(\mathcal O_{x''_\beta})
=\operatorname{long}(\mathcal O_{x'_\alpha}).
\]

